% from ega3-7-fr.tex
\oldpage[III]{219}
tout $x\in X$, si $\mathscr L$ est un $\mathscr O_X$-Module \emph{de type fini}
tel qu'il existe un $\mathscr O_X$-Module $\mathscr F$ pour lequel
$\mathscr L\otimes_{\mathscr O_X}\mathscr F$ soit \emph{isomorphe à}
$\mathscr O_X$, le raisonnement précédent montre que pour tout $x\in X$,
$\mathscr L_x$ est isomorphe à $\mathscr O_x$. On en déduit que $\mathscr L$ est
alors \emph{inversible}. En effet, pour tout $x\in X$, soient $U$ un voisinage
ouvert de $x$ tel que $\mathscr L|U$ soit engendré par $n$ sections
$s_i\ (1\leq i\leq n)$ au-dessus de $U$ (5.2.1). On peut supposer par exemple
que $s=s_1$ est telle que $s_x\ne0$, et comme $\mathscr L_x$ est isomorphe à
$\mathscr O_x$, il existe pour chaque $i$ une section $t_i$ de $\mathscr O_X$
au-dessus d'un voisinage ouvert $V_i\subset U$ de $x$ telle que
$(t_i)_x s_x=(s_i)_x$. Il y a par suite un voisinage ouvert
$V\subset\bigcap_{i=1}^nV_i$ de $x$ tel que $t_i s=s_i$ dans $V$, autrement
dit $\mathscr L|V$ est engendré par \emph{l'unique} section $s|V$. En outre,
si $z$ est une section au-dessus d'un ouvert $W\subset V$ du noyau de
l'homomorphisme $\mathscr O_X|V\to\mathscr L|V$ défini par $s$ (5.1.1), $z_y$
annule $\mathscr L_y$ pour tout $y\in W$, donc $z_y=0$ par hypothèse et par
suite $z=0$, ce qui achève de prouver notre assertion. De plus, la
considération du produit tensoriel
$\mathscr L^{-1}\otimes\mathscr L\otimes\mathscr F$ montre aussitôt que
$\mathscr F$ est isomorphe à $\mathscr L^{-1}$.

\noindent\textbf{$(\mathbf{Err}_{\mathrm{III}},\,3)$} Dans
$(0_{\mathrm I},\,7.2.4)$, il faut imposer la condition que les puissances
$\mathfrak J^n$ sont des idéaux \emph{fermés} dans l'anneau admissible $A$ pour
que la proposition soit exacte. La démonstration donnée est incorrecte, et
l'énoncé rectifié résulte de Bourbaki, \emph{Top. gén.}, chap. III,
3\textsuperscript{e} éd., \S~3, n\textsuperscript{o}~5, cor. 1 de la prop. 9.
Il faut de même supposer les $\mathfrak J^n$ \emph{fermés} dans $A$ dans les
énoncés $(0_{\mathrm I},\,7.2.5)$ et $(0_{\mathrm I},\,7.2.6)$.

\noindent\textbf{$(\mathbf{Err}_{\mathrm{III}},\,4)$} Après la proposition
(I, 2.2.5), ajouter : avec les notations de (I, 2.2.5), si $\mathfrak J$ est un
idéal de $A$, on note $\widetilde{\mathfrak J\mathscr F}$ le sous-$\mathscr O_X$-Module
$\mathfrak J\widetilde{\mathscr F}$ défini dans $(0,\,4.3.5)$.

\noindent\textbf{$(\mathbf{Err}_{\mathrm{III}},\,5)$} Dans (I, 3.4.5), ligne 1
du bas de la p. 112, remplacer « \emph{un corps} $K$ » par « \emph{un corps
algébriquement clos} $K$ ». Ligne 1 de la p. 113, remplacer « \emph{extension} »
par « \emph{extension algébriquement close} ». Lignes 1 à 4 de la p. 113,
remplacer depuis « $K$ sera appelé\ldots{} » par le texte suivant :

Pour tout point de $X$ à valeurs dans un corps $K$, $K$ sera appelé le
\emph{corps des valeurs} du point correspondant, et si $x$ est la localité de
ce point on dit encore que ce dernier est \emph{localisé en} $x$. On définit
ainsi une application $X(K)\to X$, faisant correspondre à un point à valeurs
dans $K$ sa localité.

Lignes 7 et 16 de la p. 113, supprimer « géométrique » ; lignes 10 et 13 de
la p. 113, supprimer « géométriques ». Lignes 10 et 11 du bas de la p. 115,
supprimer « \emph{géométrique} ».

\noindent\textbf{$(\mathbf{Err}_{\mathrm{III}},\,6)$} À la fin de la
démonstration de (I, 3.6.1), ligne 12 du bas de la p. 117, ajouter : Si l'on
pose
$X'=X\times_Y\operatorname{Spec}(\mathscr O_y/\mathfrak a_y)$, alors, pour tout
point $x\in X'$, identifié par $p$ à un point de $X$, on a
$\mathscr O_{X',x}=\mathscr O_{X,x}/\mathfrak a_y\mathscr O_{X,x}$. La question
étant en effet locale sur $X$ et $Y$, on peut supposer que
$X=\operatorname{Spec}(B)$, $Y=\operatorname{Spec}(A)$, et l'on a
$(B/\mathfrak a_yB)_x=B_x/\mathfrak a_yB_x$ par platitude $(0,\,1.3.2)$.

\noindent\textbf{$(\mathbf{Err}_{\mathrm{III}},\,7)$} Dans (I, 5.1.1), ligne 3
du bas de la p. 127, remplacer « $\mathscr O_X$-Module quasi-cohérent » par
« \emph{Idéal quasi-cohérent de} $\mathscr B$ ».

\noindent\textbf{$(\mathbf{Err}_{\mathrm{III}},\,8)$} Après (I, 5.1.10),
ligne 16 de la p. 131, ajouter :

\emph{Remarque (5.1.11).} --- Une légère adaptation du raisonnement fait dans
(5.1.9) montre que la conclusion reste valable \emph{sans supposer a priori
que $X$ soit un préschéma}, mais

